\section{Jordan 曲线定理}\label{ch10:section-jordan-curve}

在证明 Jordan 分割定理时所使用的 Seifert-van Kampen 定理的特殊情况告诉我们关于空间 $X=U\cup V$ 的基本群的一些信息, 其中 $U\cap V$ 是道路连通的. 在下面的定理中, 我们将考察当 $U\cap V$ 不是道路连通的时候, 情况又会怎样. 借助这个结论, 我们便可以完成 Jordan 曲线定理的证明.

\begin{theorem}\label{ch10:thm-two-open-sets-fund-group}
设 $X$ 为两个开集 $U$ 和 $V$ 之并, 并且 $U\cap V$ 可以表示成两个无交的开集 $A$ 和 $B$ 之并. 假设有一条 $U$ 中的道路 $\alpha$ 从 $A$ 的一个点 $a$ 到 $B$ 的一个点 $b$, 并有一条 $V$ 中的道路 $\beta$ 从 $b$ 到 $a$. 记 $f=\alpha*\beta$, 它是一条回路.

\begin{enumerate}[(a)]
\item 道路同伦类 $[f]$ 生成 $\pi_{1}(X,a)$ 的一个无限循环子群.
\item 若 $\pi_{1}(X,a)$ 是无限循环的, 则它是由 $[f]$ 生成的.
\item 设存在 $U$ 中的道路 $\gamma$ 从 $a$ 到 $A$ 中的点 $a'$, 又存在 $V$ 中的道路 $\delta$ 从 $a'$ 到 $a$. 于是 $g=\gamma*\delta$ 是一条回路. 这时分别由 $[f]$ 和 $[g]$ 生成的 $\pi_{1}(X,a)$ 的两个子群只交于单位元.
\end{enumerate}
\end{theorem}

\begin{proof}
这个证明中的很多地方是仿照第 54 节中证明圆周基本群是无限循环群时所采用的方法. 如前述证明一样, 这里证明的关键一步就是找到空间 $X$ 的一个适当的覆叠空间 $E$.

\textbf{第一步 (构造 $E$).} 我们通过黏贴子空间 $U$ 和 $V$ 的副本的方法构造 $E$. 选取 $U$ 的可数多个两两无交的副本和 $V$ 的可数多个两两无交的副本, 记为
\[
U\times(2n)\quad\text{和}\quad V\times(2n+1),
\]
其中 $n\in\mathbb{Z}$, $\mathbb{Z}$ 表示整数集合. 用 $Y$ 表示这些空间的并, 那么 $Y$ 便是 $X\times\mathbb{Z}$ 的一个子空间. 现在我们通过等点
\[
x\times(2n)\quad\text{和}\quad x\times(2n-1),\quad\text{对于任意 }x\in A
\]
以及等点
\[
x\times(2n)\quad\text{和}\quad x\times(2n+1),\quad\text{对于任意 }x\in B
\]
来将 $E$ 定义为 $Y$ 的商空间. 令 $\pi\colon Y\to E$ 表示商映射.

设映射 $\rho\colon Y\to X$ 定义为 $\rho(x\times m)=x$, 它诱导出映射 $p\colon E\to X$. 由于 $E$ 具有商拓扑, 所以映射 $p$ 是连续的. 另外 $p$ 也是一个满射. 下面我们来证明 $p$ 是一个覆叠映射. 参见 \cref{ch10:fig-63-1}.
\end{proof}

\par
\begin{figure}[htbp]
\centering
\includegraphics[width=0.75\textwidth]{ch10/fig-63-1.png}
\caption{定理 63.1 证明中的覆叠空间 $E$.}\label{ch10:fig-63-1}
\end{figure}

\begin{proof}[continued]
首先, 我们指出 $\pi$ 是开映射. 因为 $Y$ 是无交开集 $\{U\times(2n)\}$ 和 $\{V\times(2n+1)\}$ 的并, 因此只要证明 $\pi|(U\times 2n)$ 和 $\pi|(V\times(2n+1))$ 是开映射即可. 例如, 在 $U\times 2n$ 中任意选取一个开集, 它可以表示为 $W\times 2n$, 其中 $W$ 是 $U$ 中的一个开集. 这时
\[
\pi^{-1}(\pi(W\times 2n))=[W\times 2n]\cup[(W\cap B)\times(2n+1)]\cup[(W\cap A)\times(2n-1)]
\]
是 $Y$ 中的三个开集的并, 因而是 $Y$ 中的开集. 根据商拓扑的定义, $\pi(W\times 2n)$ 是 $E$ 中的开集.

现在我们来证明 $p$ 是一个覆叠映射. 要证明的是 $U$ 和 $V$ 能够被 $p$ 均衡地覆盖. 例如, 考虑 $U$. 集合 $p^{-1}(U)$ 是无交集合 $\pi(U\times 2n)$ ($n\in\mathbb{Z}$) 的并. 因为 $\pi$ 是开映射, 所以每一个集合都是 $E$ 中的开集. 设 $\pi_{2n}$ 表示 $\pi$ 在开集 $U\times 2n$ 上的限制, 将其映射到 $\pi(U\times 2n)$ 上. 这是一个同胚, 因为它是一一的连续开映射. 当限制到 $\pi(U\times 2n)$ 上时, 映射 $p$ 恰是两个同胚的复合
\[
\pi(U\times 2n)\xrightarrow{\pi_{2n}^{-1}}U\times 2n\xrightarrow{\rho}U,
\]
因此它也是一个同胚. 所以 $p|\pi(U\times 2n)$ 就将集合 $\pi(U\times 2n)$ 同胚地映射到 $U$ 上. 这便是我们要证明的.

\textbf{第二步.} 我们来定义回路 $f=\alpha*\beta$ 的提升的一个族.

对于每一个整数 $n$, 记 $e_{n}=\pi(a\times 2n)$ 是 $E$ 中的点. 这时诸 $e_{n}$ 是两两不同的, 并且构成集合 $p^{-1}(a)$. 定义 $f$ 的提升 $\widetilde{f}_{n}$, 它的起点为 $e_{n}$ 终点为 $e_{n+1}$.

由于 $\alpha$ 和 $\beta$ 分别是 $U$ 和 $V$ 中的道路, 我们可以定义
\[
\widetilde{\alpha}_{n}(s)=\pi(\alpha(s)\times 2n),\qquad
\widetilde{\beta}_{n}(s)=\pi(\beta(s)\times(2n+1)),
\]
这时 $\widetilde{\alpha}_{n}$ 和 $\widetilde{\beta}_{n}$ 分别是 $\alpha$ 和 $\beta$ 的提升. (当 $n=0$ 时, 如 \cref{ch10:fig-63-1} 所示.) 乘积 $\widetilde{\alpha}_{n}*\widetilde{\beta}_{n}$ 也是有定义的, 因为 $\widetilde{\alpha}_{n}$ 的终点为 $\pi(b\times 2n)$, $\widetilde{\beta}_{n}$ 的起点为 $\pi(b\times(2n+1))$. 令 $\widetilde{f}_{n}=\widetilde{\alpha}_{n}*\widetilde{\beta}_{n}$. 注意 $\widetilde{f}_{n}$ 的起点为 $\widetilde{\alpha}_{n}(0)=\pi(a\times 2n)=e_{n}$, 终点为 $\widetilde{\beta}_{n}(1)=\pi(a\times(2n+1))=\pi(a\times(2n+2))=e_{n+1}$.

\textbf{第三步.} 证明 $[f]$ 生成了 $\pi_{1}(X,a)$ 的一个无限循环子群. 也就是要证明若 $m$ 是一个正整数, 则 $[f]^{m}$ 不是单位元. 这一点是容易证明的. 因为乘积
\[
\widetilde{h}=\widetilde{f}_{0}*(\widetilde{f}_{1}*(\cdots*\widetilde{f}_{m-1}))
\]
是有定义的, 并且是 $m$-重幂
\[
h=f*(f*(\cdots*f))
\]
的提升. 这时因为 $\widetilde{h}$ 的起点为 $e_{0}$ 终点为 $e_{m}$, 所以 $[h]=[f]^{m}$ 不是平凡的.

\textbf{第四步.} 证明当 $\pi_{1}(X,a)$ 是无限循环群时, 它由 $[f]$ 生成. 考虑提升对应 $\phi\colon\pi_{1}(X,a)\to p^{-1}(a)$. 在第三步中证明了对于每一个正整数 $m$, 对应 $\phi$ 将 $[f]^{m}$ 映射为 $p^{-1}(a)$ 的一个点 $e_{m}$. 与之相似的讨论可以证明它将 $[f]^{-m}$ 映射为 $e_{-m}$. 因此 $\phi$ 是满射. 根据覆叠空间的基本定理可见, $\phi$ 诱导一个单射
\[
\Phi\colon\pi_{1}(X,a)/H\to p^{-1}(a),
\]
其中 $H=p_{*}(\pi_{1}(E,e_{0}))$. 由于 $\phi$ 是满射, 所以 $\Phi$ 是满射. 于是 $H$ 是平凡群, 因为一个无限循环群对于任何一个非平凡子群的商群是有限的. 从而提升对应 $\phi$ 是一一的. 由于这个对应将由 $[f]$ 生成的子群映射到 $p^{-1}(a)$ 上, 所以这个子群等于所有的 $\pi_{1}(X,a)$.

\textbf{第五步.} 现在我们来证明 (c) 成立. \cref{ch10:fig-63-1} 中可能误导你, 使你以为 (c) 中所考虑的 $\pi_{1}(X,a)$ 中的元 $[g]$ 是平凡的, 但实际上这是很特殊的情况. \cref{ch10:fig-63-2} 说明当 $A$ 本身是两个无交的非空开集的并时会出现什么样的情况. 在这种情况下 (稍后这对我们是有用的) $[f]$ 和 $[g]$ 都生成 $\pi_{1}(X,a)$ 的无限循环子群.
\end{proof}

\par
\begin{figure}[htbp]
\centering
\includegraphics[width=0.75\textwidth]{ch10/fig-63-2.png}
\caption{当 $A$ 由两个无交开集组成时的提升.}\label{ch10:fig-63-2}
\end{figure}

\begin{proof}[continued]
给定 $g=\gamma*\delta$, 我们定义 $g$ 到 $E$ 上的一个提升如下: 因为 $\gamma$ 是 $U$ 中的一条道路, 我们可以定义
\[
\widetilde{\gamma}(s)=\pi(\gamma(s)\times 0);
\]
因为 $\delta$ 是 $V$ 中的一条道路, 我们可以定义
\[
\widetilde{\delta}(s)=\pi(\delta(s)\times(-1)).
\]
这时 $\widetilde{\gamma}$ 和 $\widetilde{\delta}$ 分别是 $\gamma$ 和 $\delta$ 的提升. 乘积 $\widetilde{g}=\widetilde{\gamma}*\widetilde{\delta}$ 有定义, 因为 $\widetilde{\gamma}$ 的终点为 $\pi(a'\times 0)$, $\widetilde{\delta}$ 的起点为 $\pi(a'\times(-1))$. 此外它还是 $g$ 的提升. 注意 $\widetilde{g}$ 是 $E$ 中的一条回路, 因为它的起点和终点都是 $\pi(a\times 0)=\pi(a\times(-1))=e_{0}$.

这蕴涵着由 $[f]$ 和 $[g]$ 生成的两个子群的公共部分只有单位元. 因为 $f$ 的 $m$-重幂提升到一条起点是 $e_{0}$ 终点是 $e_{m}$ 的道路, 但是 $g$ 的任意重幂提升到一条起点和终点都是 $e_{0}$ 的道路, 所以对于每一个非零的整数 $m$ 和 $k$, 有 $[f]^{m}\neq[g]^{k}$.
\end{proof}

\begin{theorem}[nonseparation theorem]\label{ch10:thm-nonseparation}
设 $D$ 是 $S^{2}$ 中的一段弧. 则 $D$ 不分割 $S^{2}$.
\end{theorem}

\begin{proof}
我们给出这个定理的两个证明. 第一个证明使用了前一节中的结论, 第二个证明则没有.

\textbf{第一个证明.} 由于 $D$ 是可缩的, 所以恒等映射 $i\colon D\to D$ 是零伦的. 因此若 $a$ 和 $b$ 是 $S^{2}$ 中的不属于 $D$ 的两个点, 则内射 $j\colon D\to S^{2}-a-b$ 是零伦的. 由 Borsuk 引理可见 $a,b$ 属于 $S^{2}-D$ 的同一分支.

\textbf{第二个证明.} 将 $D$ 表示为两段弧 $D_{1}$ 和 $D_{2}$ 的并, 并且这两段弧只交于一点 $d$. 设 $a$ 和 $b$ 是 $S^{2}$ 中的不属于 $D$ 的两个点. 我们证明若 $a,b$ 能被 $S^{2}-D_{1}$ 和 $S^{2}-D_{2}$ 中的道路连起来, 那么它们就能被 $S^{2}-D$ 中的一条道路连接. \cref{ch10:fig-63-3} 说明了这一论断不是平凡的.
\end{proof}

\par
\begin{figure}[htbp]
\centering
\includegraphics[width=0.5\textwidth]{ch10/fig-63-3.png}
\caption{弧 $D=D_{1}\cup D_{2}$ 交于 $d$.}\label{ch10:fig-63-3}
\end{figure}

\begin{proof}[continued]
现在假设 $a$ 和 $b$ 不能被 $S^{2}-D$ 中的道路连接, 并且由此推出矛盾. 应用 \cref{ch10:thm-two-open-sets-fund-group}. 令 $X$ 表示空间 $S^{2}-d$. 令 $U$ 和 $V$ 分别表示两个开集
\[
U=S^{2}-D_{1}\quad\text{和}\quad V=S^{2}-D_{2},
\]
则 $X=U\cup V$, 并且 $U\cap V=S^{2}-D$. 根据假设, $a,b$ 是 $S^{2}-D$ 中的两个不能由 $S^{2}-D$ 中的道路连接的点. 因此, $U\cap V$ 不是道路连通的. 设 $A$ 是 $U\cap V$ 中的包含 $a$ 的道路分支, $B$ 是其他道路分支的并. 因为 $U\cap V$ (作为 $S^{2}$ 中的开集) 是局部道路连通的, 所以 $U\cap V$ 的道路分支是开集, 从而 $A$ 和 $B$ 都是 $X$ 中的开集. 根据假定, $a$ 和 $b$ 既可以被 $U=S^{2}-D_{1}$ 中的道路连接, 也可以被 $V=S^{2}-D_{2}$ 中的道路连接. 根据 \cref{ch10:thm-two-open-sets-fund-group} 有 $\pi_{1}(X,a)$ 不是平凡的. 然而 $X=S^{2}-d$, 因此它的基本群却是平凡的.

现在来完成定理的证明. 给定弧 $D$ 以及 $S^{2}-D$ 中的两个点 $a$ 和 $b$, 假设 $a$ 和 $b$ 不能被 $S^{2}-D$ 中的道路连接, 并且由此推出矛盾. 选取同胚 $h\colon[0,1]\to D$, 令 $D_{1}=h([0,1/2])$, $D_{2}=h([1/2,1])$. 上面已经说明了因为 $a$ 和 $b$ 不能被 $S^{2}-D$ 中的道路连接, 所以它们既不能被 $S^{2}-D_{1}$ 中的道路连接, 又不能被 $S^{2}-D_{2}$ 中的道路连接. 不妨设 $a,b$ 不能被 $S^{2}-D_{1}$ 中的道路连接.

将 $D_{1}$ 分割为两段弧 $E_{1}=h([0,1/4])$ 和 $E_{2}=h([1/4,1/2])$, 然后重复前面的论证, 我们便可证明 $a$ 和 $b$ 不能同时在 $S^{2}-E_{1}$ 和 $S^{2}-E_{2}$ 中用道路连接.

类似地, 不断地重复这个论证过程. 我们便可定义闭区间的一个序列
\[
I\supset I_{1}\supset I_{2}\supset\cdots
\]
使得 $I_{n}$ 的长度为 $(1/2)^{n}$, 并且对任意 $n$, $a$ 和 $b$ 不能被 $S^{2}-h(I_{n})$ 中的道路连接. 根据单位区间的紧致性, 存在 $x$ 属于 $\bigcap I_{n}$. 又因区间的长度趋于零, 所以这样的点只有一个.

考虑空间 $S^{2}-h(x)$. 因为这个空间同胚于 $\mathbb{R}^{2}$, 所以点 $a$ 和 $b$ 能够被 $S^{2}-h(x)$ 中的道路 $\alpha$ 连接. 由于 $\alpha(I)$ 是紧致的, 因此是闭的. 从而存在 $h(x)$ 的一个 $\varepsilon$-邻域与 $\alpha(I)$ 无交. 又因为 $h$ 是连续的, 所以存在 $m$ 使得 $h(I_{m})$ 包含在这个邻域中. 这就说明 $\alpha$ 是 $S^{2}-h(I_{m})$ 中连接 $a,b$ 的道路, 与假设矛盾.
\end{proof}

这个定理的两个证明都很有趣. 正如我们在第 62 节中注意到的那样, 第一个证明是为了说明不存在 $S^{2}$ 的紧致可缩子空间能分割 $S^{2}$. 第二个证明是从另一个方面入手的. 我们来检查一下第二个证明, 并且考虑一下到底 $D_{1}$ 和 $D_{2}$ 的哪些性质使得结论成立. 事实上, 只需要 $D_{1}$ 和 $D_{2}$ 是 $S^{2}$ 的闭子集并且 $S^{2}-(D_{1}\cap D_{2})$ 是单连通的. 因此下面的定理也成立, 稍后我们还会用到这个定理.

\begin{theorem}[general nonseparation theorem]\label{ch10:thm-general-nonseparation}
设 $D_{1}$ 和 $D_{2}$ 是 $S^{2}$ 的闭子集并且 $S^{2}-D_{1}\cap D_{2}$ 是单连通的. 若 $D_{1}$ 和 $D_{2}$ 都不分割 $S^{2}$, 则 $D_{1}\cup D_{2}$ 不分割 $S^{2}$.
\end{theorem}

现在我们来证明 Jordan 曲线定理.

\begin{theorem}[Jordan curve theorem]\label{ch10:thm-jordan-curve}
设 $C$ 是 $S^{2}$ 中的一条简单闭曲线, 则 $C$ 恰好将 $S^{2}$ 分割成两个分支 $W_{1}$ 和 $W_{2}$, 并且 $W_{1}$ 和 $W_{2}$ 都将 $C$ 作为它的边界, 即 $C=\overline{W}_{i}-W_{i}$, $i=1,2$.
\end{theorem}

\begin{proof}
\textbf{第一步.} 证明 $S^{2}-C$ 恰有两个分支. 将 $C$ 表示为两段弧 $C_{1}$ 和 $C_{2}$ 的并, 其中 $C_{1}$ 和 $C_{2}$ 交于两点集 $\{p,q\}$. 令 $X$ 为空间 $S^{2}-p-q$, $U$ 和 $V$ 为两个开集
\[
U=S^{2}-C_{1}\quad\text{和}\quad V=S^{2}-C_{2},
\]
这时 $X=U\cup V$, 并且 $U\cap V=S^{2}-C$. 根据 Jordan 分割定理, $U\cap V$ 至少有两个分支.

假定 $U\cap V$ 的分支多于两个, 并且由此推出矛盾. 设 $A_{1}$ 和 $A_{2}$ 是 $U\cap V$ 的两个分支, $B$ 是其他分支的并. 因为 $S^{2}-C$ 是局部连通的, 所以这些集合都是开集. 设 $a\in A_{1}$, $a'\in A_{2}$, $b\in B$. 因为弧 $C_{1}$ 和 $C_{2}$ 不能分割 $S^{2}$, 所以有 $U$ 中的道路 $\alpha$ 和 $\gamma$ 分别从 $a$ 到 $b$ 和 $a$ 到 $a'$, 也有 $V$ 中的道路 $\beta$ 和 $\delta$ 分别从 $b$ 到 $a$ 和 $a'$ 到 $a$. 考虑回路 $f=\alpha*\beta$ 和 $g=\gamma*\delta$. 将 $U\cap V$ 表示成两个开集 $A_{1}\cup A_{2}$ 和 $B$ 的并, 根据 \cref{ch10:thm-two-open-sets-fund-group}, $[f]$ 是 $\pi_{1}(X,a)$ 中的非平凡元. 将 $U\cap V$ 表示成两个开集 $A_{1}$ 和 $A_{2}\cup B$ 的并, 那么 $[g]$ 也是 $\pi_{1}(X,a)$ 中的非平凡元. 因为 $\pi_{1}(X,a)$ 是无限循环的, 所以存在非零整数 $m,k$ 使得 $[f]^{m}=[g]^{k}$, 但这与 \cref{ch10:thm-two-open-sets-fund-group} 的 (c) 矛盾.

\textbf{第二步.} 证明 $C$ 是 $W_{1}$ 和 $W_{2}$ 的公共边界.

因为 $S^{2}$ 是局部连通的, 所以 $S^{2}-C$ 的两个分支 $W_{1}$ 和 $W_{2}$ 都是 $S^{2}$ 中的开集. 特别地, 两者都不包含对方的极限点, 从而 $\overline{W}_{1}-W_{1}$ 和 $\overline{W}_{2}-W_{2}$ 都包含在 $C$ 中.

为了证明反向的包含关系, 我们需要证明 $C$ 中的每一个点 $x$ 的任何一个邻域 $U$ 都与闭集 $\overline{W}_{1}-W_{1}$ 有非空的交, 从而说明 $x$ 属于 $\overline{W}_{1}-W_{1}$.

设 $U$ 是 $x$ 的任意一个邻域. 因为 $C$ 同胚于 $S^{1}$, 我们可将 $C$ 表示成只交于端点的两段弧 $C_{1}$ 和 $C_{2}$, 并且 $C_{1}$ 充分小, 使得它能包含在 $U$ 中. 参见 \cref{ch10:fig-63-4}.

设 $a$ 和 $b$ 分别是 $W_{1}$ 和 $W_{2}$ 中的点. 因为 $C_{2}$ 不分割 $S^{2}$, 所以我们可以找到 $S^{2}-C_{2}$ 中的一条道路 $\alpha$ 连接 $a$ 和 $b$. 集合 $\alpha(I)$ 必定包含 $\overline{W}_{1}-W_{1}$ 的一个点 $y$, 否则 $\alpha(I)$ 就是包含在无交开集 $W_{1}$ 和 $S^{2}-\overline{W}_{1}$ 的并中的连通集, 并且还会和两个开集都相交. 点 $y$ 属于闭曲线 $C$, 因为 $(\overline{W}_{1}-W_{1})\subset C$. 因为道路 $\alpha$ 与弧 $C_{2}$ 无交, 所以 $y$ 必包含于弧 $C_{1}$, 因此便属于开集 $U$. 从而 $U$ 与 $\overline{W}_{1}-W_{1}$ 交于点 $y$. 定理证毕.
\end{proof}

\par
\begin{figure}[htbp]
\centering
\includegraphics[width=0.5\textwidth]{ch10/fig-63-4.png}
\caption{$C=\overline{W}_{1}-W_{1}=\overline{W}_{2}-W_{2}$.}\label{ch10:fig-63-4}
\end{figure}

如同我们对前述定理所做的一样, 我们想知道使定理成立的关键条件到底是什么. 回头细看第一步, 实际上只需要 $C_{1}$ 和 $C_{2}$ 是连通闭集, $C_{1}\cap C_{2}$ 只有两个点, 以及 $C_{1}$ 和 $C_{2}$ 都不分割 $S^{2}$. 前两个条件保证了 $C_{1}\cup C_{2}$ 至少能够分割 $S^{2}$ 为两个分支, 第三个条件保证了它只能被分割为两个分支. 从而, 我们不需作更多的努力就可以获得下面的结论.

\begin{theorem}\label{ch10:thm-two-arcs-nonseparating}
设 $C_{1}$ 和 $C_{2}$ 是 $S^{2}$ 中的仅交于两点的两个连通子集. 如果 $C_{1}$ 和 $C_{2}$ 都不分割 $S^{2}$, 那么 $C_{1}\cup C_{2}$ 将 $S^{2}$ 分割成两个分支.
\end{theorem}

\begin{example}\label{ch10:ex-jordan-curve-1}
读者可能会觉得 Jordan 曲线定理的第二部分 (即 $C$ 是 $W_{1}$ 和 $W_{2}$ 的公共边界) 太明显了用不着证明. 其实这个结论成立严格依赖于 $C$ 同胚于 $S^{1}$.

例如, 考虑 \cref{ch10/fig-63-5} 中的空间. 它是相交于两点的两段弧的并, 根据 \cref{ch10:thm-two-arcs-nonseparating}, 它将 $S^{2}$ 分割为两个分支 $W_{1}$ 和 $W_{2}$, 就像圆周所做的一样. 然而在这种情形下 $C$ 并不是 $W_{1}$ 和 $W_{2}$ 的公共边界.
\end{example}

\par
\begin{figure}[htbp]
\centering
\includegraphics[width=0.6\textwidth]{ch10/fig-63-5.png}
\caption{两段弧之并分割 $S^{2}$, 但不是两分支的公共边界.}\label{ch10/fig-63-5}
\end{figure}

除了这三个分割定理外还有第四个分割定理, 称为 Schoenflies 定理. 这个定理说: 若 $C$ 是 $S^{2}$ 中的一条简单闭曲线, $U$ 和 $V$ 是 $S^{2}-C$ 的两个分支, 则 $\overline{U}$ 和 $\overline{V}$ 同胚于单位闭球 $B^{2}$. 这个定理的一个证明可以在 $[\mathrm{H-S}]$ 中找到.

分割定理可以推广到高维情形, 如下:

\begin{enumerate}[(1)]
\item $S^{n}$ 中任意同胚于 $S^{n-1}$ 的子空间 $C$ 分割 $S^{n}$.
\item $S^{n}$ 中任意同胚于 $[0,1]$ 或某一个球 $B^{m}$ 的子空间 $A$ 不分割 $S^{n}$.
\item $S^{n}$ 中任意同胚于 $S^{n-1}$ 的子空间 $C$ 分割 $S^{n}$ 为两个分支, $C$ 是这两个分支的公共边界.
\end{enumerate}

学习了代数拓扑中的奇异同调群就可以容易地证明这些定理. (参见 $[\mathrm{Mu}]$, 202 页.) 关于 $\mathbb{R}^{n}$ 的 Brouwer 区域不变性定理只是一个推论.

然而, 如果不对空间 $C$ 到 $S^{n}$ 的嵌入加一些限制条件, Schoenflies 定理便不能推广到高维. 著名的例子``Alexander 角球''便说明了这一点. Alexander 角球是 $S^{2}$ 在 $S^{3}$ 中的一个同胚像, 它的补区域中竟然有一个不是单连通的. (参见 $[\mathrm{H-Y}]$, 176 页.)

这些分割定理能够更进一步推广. 按照这条路线走到最后便是著名的 Alexander-Pontryagin 对偶定理, 这是一个极深刻的代数拓扑定理, 在这里我们不打算陈述它. (参见 $[\mathrm{Mu}]$.) 从 Alexander-Pontryagin 对偶定理可以推出: 如果 $S^{n}$ 的某一个闭子空间 $C$ 将 $S^{n}$ 分割为 $k$ 个分支, 那么 $S^{n}$ 的每一个同胚于 (甚至是同伦等价于) $C$ 的子空间也会将 $S^{n}$ 分割为 $k$ 个分支. 这样一来, 前面提到的高维情形下的三个定理 (1)~(3) 便可以立即由此推出.

\begin{exercise}\label{ch10:ex-jordan-curve-1}
设 $C_{1}$ 和 $C_{2}$ 是 $S^{2}$ 中的两条无交的简单闭曲线.

\begin{enumerate}[(a)]
\item 证明 $S^{2}-C_{1}-C_{2}$ 只有三个分支. [提示: 若 $W_{1}$ 是 $S^{2}-C_{1}$ 的分支, 与 $C_{2}$ 无交, 并且 $W_{2}$ 是 $S^{2}-C_{2}$ 的分支, 与 $C_{1}$ 无交, 证明 $W_{1}\cup\overline{W}_{2}$ 不能分割 $S^{2}$.]
\item 证明这三个分支的边界分别为 $C_{1}$, $C_{2}$ 和 $C_{1}\cup C_{2}$.
\end{enumerate}
\end{exercise}

\begin{exercise}\label{ch10:ex-jordan-curve-2}
设 $D$ 是 $S^{2}$ 的闭连通子空间, 它将 $S^{2}$ 分割成 $n$ 个分支.

\begin{enumerate}[(a)]
\item 若 $A$ 是 $S^{2}$ 中的一段弧, 与 $D$ 的交为它的一个端点, 证明 $D\cup A$ 将 $S^{2}$ 分割成 $n$ 个分支.
\item 若 $A$ 是 $S^{2}$ 中的一段弧, 与 $D$ 的交为它的两个端点, 证明 $D\cup A$ 将 $S^{2}$ 分割成 $n+1$ 个分支.
\item 若 $C$ 是 $S^{2}$ 中的一条简单闭曲线, 与 $D$ 仅交于一点, 证明 $D\cup C$ 将 $S^{2}$ 分割成 $n+1$ 个分支.
\end{enumerate}
\end{exercise}

\begin{exercise}\label{ch10:ex-jordan-curve-3}
设 $D$ 是 $S^{2}$ 的子空间, 同胚于拓扑学家的正弦曲线 $\overline{S}$. (参见第 24 节.) 证明 $D$ 不分割 $S^{2}$. [提示: 设 $h\colon\overline{S}\to D$ 是一个同胚. 给定 $0<c<1$, 令 $\overline{S}_{c}$ 表示 $\overline{S}$ 与集合 $\{(x,y)\mid x\leqslant c\}$ 的交. 证明: 对于给定的 $a,b\in S^{2}-D$, 对某一个 $c$, 存在 $S^{2}-h(\overline{S}_{c})$ 中的一条从 $a$ 到 $b$ 的道路. 由此推出存在 $S^{2}-D$ 中的一条从 $a$ 到 $b$ 的道路.]
\end{exercise}
\begin{enumerate}[(b)]
\item 设 $C$ 是 $S^{2}$ 的子空间, 同胚于闭的拓扑学家的正弦曲线. 证明 $C$ 正好将 $S^{2}$ 分割成两个分支, 并且 $C$ 是这两个分支的公共边界. [提示: 设 $h$ 是闭的拓扑学家的正弦曲线与 $C$ 之间的一个同胚. 设 $C_{0}=h(0\times[-1,1])$. 先应用 Jordan 曲线定理的论断证明 $C-C_{0}$ 中的每一个点都是 $S^{2}-C$ 的每一个分支的边界点.]
\end{enumerate}
